% ============================================================================
% Fichier  : partie4/P04_C25_audit_conformite.tex
% Titre    : Auditer la Conformité
% Auteur   : MINKA MI NGUIDJOÏ Thierry Emmanuel
% Version  : 1.0 -- novembre 2026
% ============================================================================
% Base : ACSI C12 (paiements de l'ANPP, 40 = 29/6/5, µ 0,73 U 0,13 H 1,12 ;
% 4 agents payés après départ pour 1 424 000 FCFA, 1 ordonnateur =
% comptable, 1 paiement sans visa ; 5 pièces « aux archives ») et ACSI C10
% (relation tripartite, ISSAI 400.35 ; règle de référence ISSAI 4000.14).
% Sources : INTOSAI ISSAI 100, 400, 4000 ; IDI, manuel d'audit de
% conformité (2021) ; ISO 37301:2021 (systèmes de management de la
% conformité) ; loi n° 2024/017 (données personnelles, présentée sans
% citation, à vérifier au Journal officiel) ; régime financier de l'État.
% Corrigés : corriges/Corr_C25.tex.
% ============================================================================

\chapter{Auditer la Conformité}
\label{chap:audit-conformite}
\chaptermark{Auditer la conformité}

\epigraph{%
    « La conformité ne se présume pas.\\
    Elle se démontre, obligation par obligation. »%
}{Principe de l'audit de conformité}

\niveauChapitre{Intermédiaire}{%
    Maîtrise des Chapitres~\ref{chap:normes} (normes de l'INTOSAI) et
    \ref{chap:droit} (du texte au critère)%
}

\begin{boxobjectifs}
\begin{itemize}
    \item Définir l'audit de conformité, sa relation à trois parties et ses
          deux niveaux d'assurance.
    \item Conduire un audit de conformité selon l'ISSAI~4000, de la
          planification au rapport.
    \item Construire un registre des obligations et en tirer les critères.
    \item Formuler une conclusion ou une opinion de conformité, avec ou sans
          réserve.
    \item Auditer la conformité au droit des données personnelles.
    \item Distinguer l'audit de conformité de la fonction conformité.
\end{itemize}
\end{boxobjectifs}

\bigskip

% ============================================================================
\section*{Question directrice : l'ANPP respecte-t-elle les règles de la dépense~?}
% ============================================================================

La Chambre des comptes a relevé, dans un premier sondage, des paiements de
l'ANPP sans pièce justificative (Chapitre~\ref{chap:droit}). Le Parlement
demande une conclusion~: l'agence respecte-t-elle, oui ou non, les règles
de la dépense publique dans ses paiements de 2025~? M.~NJOYA, magistrat
vérificateur, sait qu'il doit répondre par une conclusion de conformité, et
que cette conclusion n'a de valeur que si les critères, la méthode et les
preuves sont irréprochables.

\medskip
\textbf{Plan du chapitre.}
\begin{enumerate}
    \item Ce qu'est un audit de conformité (section~\ref{sec:c25-definition}).
    \item La méthode de l'ISSAI~4000 (section~\ref{sec:c25-methode}).
    \item Le registre des obligations (section~\ref{sec:c25-registre}).
    \item Conclure sur la conformité (section~\ref{sec:c25-conclure}).
    \item Les paiements de l'ANPP (section~\ref{sec:c25-anpp}).
    \item La conformité au droit des données personnelles
          (section~\ref{sec:c25-donnees}).
\end{enumerate}

\separateur

% ============================================================================
\section{Ce qu'est un Audit de Conformité}
\label{sec:c25-definition}
% ============================================================================

\begin{boxdef}
\termecle{Audit de conformité}\index{audit de conformité}~: évaluation
indépendante de la mesure dans laquelle un objet (une activité, des
opérations, des informations) respecte, en tous points significatifs, les
critères qui lui sont applicables~: lois, règlements, contrats, règles
internes~\cite{idi2021conformite}.
\end{boxdef}

\subsection{Trois parties}

L'ISSAI~400 décrit la relation à trois parties qui fonde tout audit de
conformité~: la \emph{partie responsable} (l'entité qui doit respecter les
règles), l'\emph{auditeur}, et les \emph{utilisateurs prévus} (ceux qui se
fieront à la conclusion). Le Chapitre~\ref{chap:normes} l'a appliquée à
l'ANPP~: le directeur général M.~ESSOMBA, la Chambre des comptes, le
Parlement et le public.

\subsection{Deux niveaux d'assurance}

\begin{tableaucours}{L{3.0cm}L{4.4cm}L{4.8cm}}{%
    \tableauHeader{Assurance} &
    \tableauHeader{Travaux} &
    \tableauHeader{Forme de la conclusion}
}
Raisonnable & Tests étendus, preuves suffisantes pour une opinion & « En
tous points significatifs, les paiements sont conformes \ldots » \\
Limitée & Travaux réduits (analyses, entretiens) & « Nous n'avons rien
relevé qui nous porte à croire que les paiements ne sont pas conformes
\ldots » \\
\end{tableaucours}
\vspace{-0.8em}
\begin{center}{\small\itshape Tableau~25.1. Assurance raisonnable et
assurance limitée}\end{center}

\begin{boxinfo}
\textbf{Audit de conformité et fonction conformité.} La fonction
conformité (deuxième ligne) veille, au quotidien, à ce que l'organisation
connaisse et respecte ses obligations~; l'ISO~37301 en propose un système
de management~\cite{iso37301_2021}. L'audit de conformité (troisième ligne
ou contrôle externe) évalue, à un moment donné, le résultat. Le premier
produit le registre des obligations~; le second s'en sert.
\end{boxinfo}

% ============================================================================
\section{La Méthode de l'ISSAI 4000}
\label{sec:c25-methode}
% ============================================================================

L'ISSAI~4000 et le manuel de l'IDI déroulent l'audit de conformité en
trois phases, que la méthode du cours outille~\cite{idi2021conformite}.

\begin{tableaucours}{L{2.6cm}L{5.4cm}L{4.2cm}}{%
    \tableauHeader{Phase} &
    \tableauHeader{Étapes} &
    \tableauHeader{Outil du cours}
}
Planification & Comprendre l'entité et son contrôle interne~; identifier
l'objet et les critères~; apprécier les risques de non-conformité~; fixer
le seuil de signification et la stratégie & Lettre de mission, registre des
obligations, décomposition, fiches d'assertion \\
Exécution & Recueillir les preuves (tests sur échantillon ou sur
population)~; évaluer les écarts & Cotation \OUI{}, \NON{}, \NEK{}~; papiers
de travail \\
Rapport & Conclusion ou opinion~; constats~; réponses de la partie
responsable~; recommandations & Lecture $\mu$, $U$, $H$~; observation en
sept éléments \\
\end{tableaucours}
\vspace{-0.8em}
\begin{center}{\small\itshape Tableau~25.2. Les phases de l'audit de
conformité (d'après l'ISSAI~4000 et l'IDI)}\end{center}

\subsection{Les critères d'abord}

Un audit de conformité sans critères écrits n'est pas un audit~: c'est une
opinion. L'ISSAI~4000 exige que les critères soient identifiés avant le
terrain et communiqués à la partie responsable. Le
Chapitre~\ref{chap:droit} a montré comment transformer un texte en critères
atomiques~; ce chapitre les rassemble dans un registre.

\subsection{Le seuil de signification}

\begin{boxdef}
\termecle{Seuil de signification}\index{seuil de signification}~: niveau
au-delà duquel une non-conformité, seule ou cumulée, influencerait la
décision des utilisateurs prévus. Il peut être quantitatif (un montant) ou
qualitatif (une infraction grave est significative quel que soit son
montant).
\end{boxdef}

Pour les paiements de l'ANPP, M.~NJOYA retient deux seuils~: un seuil
quantitatif de 1~\% du total des paiements testés, et un seuil qualitatif~:
tout paiement à une personne non due (agent parti, bénéficiaire fictif) est
significatif quel qu'en soit le montant.

% ============================================================================
\section{Le Registre des Obligations}
\label{sec:c25-registre}
% ============================================================================

\begin{boxprocedure}
\textbf{Construire un registre des obligations}
\begin{enumerate}
    \item Recenser les textes applicables (lois, règlements, actes
          uniformes, contrats, règles internes), avec leur version et leur
          source.
    \item Pour chaque texte, isoler les obligations (sujet, action, objet,
          délai) et les écrire de façon atomique.
    \item Désigner le propriétaire de chaque obligation dans l'organisation.
    \item Associer à chaque obligation le contrôle qui en assure le
          respect, et la preuve attendue.
    \item Dater le registre, et le revoir à chaque nouveau texte.
\end{enumerate}
\end{boxprocedure}

\begin{tableaucours}{cL{3.8cm}L{2.4cm}L{2.6cm}L{2.6cm}}{%
    \tableauHeader{N°} &
    \tableauHeader{Obligation} &
    \tableauHeader{Source} &
    \tableauHeader{Contrôle} &
    \tableauHeader{Preuve}
}
1 & Le paiement est précédé d'un ordonnancement & Régime financier & Visa
de l'ordonnateur & Mandat daté \\
2 & Ordonnateur et comptable sont distincts & Régime financier &
Séparation des droits dans le système & Matrice des droits, mandat \\
3 & Le paiement porte le visa du contrôle financier & Régime financier &
Blocage sans visa & Mandat visé \\
4 & Le service fait est attesté & Régime financier & Pièce de service fait
& Attestation datée \\
5 & Le bénéficiaire est un agent en fonction à la date du service &
Statut des agents & Rapprochement paie et effectifs & Fichier des
départs \\
\end{tableaucours}
\vspace{-0.8em}
\begin{center}{\small\itshape Tableau~25.3. Extrait du registre des
obligations de l'ANPP pour les paiements (régime financier présenté
d'après le cours~; textes à vérifier)}\end{center}

% ============================================================================
\section{Conclure sur la Conformité}
\label{sec:c25-conclure}
% ============================================================================

\subsection{Quatre formes de conclusion}

\begin{tableaucours}{L{3.2cm}L{4.6cm}L{4.4cm}}{%
    \tableauHeader{Conclusion} &
    \tableauHeader{Quand} &
    \tableauHeader{Lecture du cours}
}
Sans réserve & Aucune non-conformité significative & Conformité claire \\
Avec réserve & Non-conformités significatives mais circonscrites, ou
limitation d'étendue partielle & Conformité claire avec $\mu$ élevé mais
des \NON{} significatifs~; ou $U$ localisé \\
Défavorable & Non-conformités significatives et généralisées &
Non-conformité claire \\
Impossibilité de conclure & Preuves insuffisantes sur une part
significative & Domaine opaque ($U$ élevé) \\
\end{tableaucours}
\vspace{-0.8em}
\begin{center}{\small\itshape Tableau~25.4. Formes de conclusion de
conformité et lecture du cours}\end{center}

\begin{boxinfo}
\textbf{\NEK{} et limitation d'étendue.} Dans le vocabulaire de l'audit de
conformité, une pièce qu'on ne peut pas obtenir est une \emph{limitation
d'étendue}. La cotation \NEK{} en est la trace pièce par pièce~; le taux
$U$ dit si la limitation est assez large pour empêcher toute conclusion.
\end{boxinfo}

% ============================================================================
\section{Les Paiements de l'ANPP}
\label{sec:c25-anpp}
% ============================================================================

\begin{boxlecture}{Paiements de l'ANPP}
M.~NJOYA examine un échantillon aléatoire de 40~paiements de l'ANPP au
regard des obligations du registre.

\medskip
\textbf{Constatations.} 29 paiements sont conformes à toutes les
obligations (\OUI{}). 6 ne le sont pas (\NON{})~: quatre ont été versés à
des agents après la date de leur départ, pour 1~424~000~FCFA~; un a le même
agent pour ordonnateur et pour comptable~; un n'a pas de visa du contrôle
financier. Pour 5 paiements, les pièces sont déclarées « aux archives » et
n'ont pas été produites (\NEK{}).

\medskip
$\mu = 29/40 \approx 0{,}73$~; $U = 5/40 \approx 0{,}13$~; $H \approx
1{,}12$~bit.

\medskip
\textbf{Analyse} (confiance élevée). Conformité claire au sens du quadrant
($\mu \ge 0{,}5$, $U < 0{,}20$), mais les quatre paiements à des agents
partis franchissent le seuil qualitatif~: la conclusion ne peut pas être
sans réserve. Les cinq \NEK{} sont une limitation d'étendue, à mentionner.

\medskip
\hfill\lecture{0{,}73}{0{,}13}{1{,}12}
\end{boxlecture}

\begin{boxscenario}{Opération MEFOU. La conclusion de M.~NJOYA}
\textit{« Sur la base des travaux réalisés sur un échantillon aléatoire de
40~paiements de l'exercice~2025, et à l'exception des points décrits
ci-après, les paiements de l'ANPP sont conformes, en tous points
significatifs, aux règles de la dépense publique. Points ayant motivé la
réserve~: quatre paiements, pour 1~424~000~FCFA, ont été versés à des agents
après la date de leur départ~; un paiement a été ordonnancé et payé par le
même agent~; un paiement n'a pas reçu le visa du contrôle financier.
Limitation~: les pièces de cinq paiements n'ont pas été produites. »}
\end{boxscenario}

\begin{boxalert}
\textbf{Violation typique.} Extrapoler à toute la population le montant des
quatre paiements irréguliers (« l'ANPP a payé 1,4~million sur 40, donc
environ X~millions sur l'année ») sans méthode statistique.

\smallskip
\textbf{Règle~:} un échantillon aléatoire permet une projection, mais
seulement avec la méthode adaptée (Chapitre~\ref{chap:preuve}) et un
intervalle. Sans elle, on rapporte les constats de l'échantillon, et on
recommande un test sur toute la population (rapprochement des paiements et
du fichier des départs).
\end{boxalert}

% ============================================================================
\section{La Conformité au Droit des Données Personnelles}
\label{sec:c25-donnees}
% ============================================================================

La loi n°~2024/017 encadre la protection des données à caractère personnel
au Cameroun~; son texte est à vérifier au Journal officiel avant toute
citation. Sans préjuger de ses articles, un audit de conformité au droit
des données personnelles porte, dans la plupart des législations de ce
type, sur un ensemble d'obligations que l'auditeur retrouve dans le texte
national avant de les écrire en critères.

\begin{tableaucours}{L{3.4cm}L{5.0cm}L{3.8cm}}{%
    \tableauHeader{Thème} &
    \tableauHeader{Question d'audit} &
    \tableauHeader{Preuve}
}
Inventaire des traitements & Les traitements de données personnelles
sont-ils recensés~? & Registre des traitements \\
Finalité et base & Chaque traitement a-t-il une finalité et un fondement~?
& Registre, mentions d'information \\
Minimisation et durée & Ne collecte-t-on que le nécessaire, pour une durée
fixée~? & Formulaires, règles de purge \\
Sécurité & Les mesures de sécurité sont-elles adaptées au risque~? &
Mesures techniques (Chapitre~\ref{chap:audit-technique}) \\
Droits des personnes & Les demandes d'accès et de rectification sont-elles
traitées~? & Registre des demandes \\
Violations & Les violations sont-elles détectées, notifiées et
documentées~? & Registre des violations \\
Transferts et sous-traitants & Les transferts et les prestataires sont-ils
encadrés~? & Contrats (Chapitre~\ref{chap:audit-tiers}) \\
\end{tableaucours}
\vspace{-0.8em}
\begin{center}{\small\itshape Tableau~25.5. Thèmes d'un audit de conformité
au droit des données personnelles (à rapprocher du texte de la loi
n°~2024/017)}\end{center}

\begin{boxscenario}{Opération LIZBIZ. Les données des clients}
\textbf{Faits.} LiZbiZ traite les données de 12~000~clients (identité,
téléphone, adresses de livraison). Il n'existe aucun registre des
traitements~; le portail de suivi présente une faille critique
(Chapitre~\ref{chap:intrusion})~; le prestataire de Garoua accède aux
serveurs sans clause de confidentialité spécifique aux données.

\medskip
\textbf{Analyse} (confiance moyenne). Sur les sept thèmes, trois sont
manifestement défaillants (inventaire, sécurité, sous-traitants)~; les
autres ne peuvent pas être jugés sans registre (\NEK{}). La première
recommandation est de constituer le registre des traitements~: sans lui,
aucune conformité ne se démontre.
\end{boxscenario}

\begin{boxmetier}{Responsable de la conformité}
\textbf{Tenir le registre des obligations}
\begin{enumerate}
    \item Recenser les textes et leur version, avec la date de lecture.
    \item Écrire chaque obligation de façon atomique, avec son
          propriétaire.
    \item Associer à chaque obligation un contrôle et une preuve.
    \item Vérifier chaque année un échantillon d'obligations, et présenter
          le résultat au comité d'audit.
    \item Mettre le registre à jour à chaque nouveau texte (par exemple, à
          la publication des textes d'application de la loi n°~2024/017).
\end{enumerate}
\end{boxmetier}

\begin{boxmetier}{Auditeur SI}
\textbf{Préparer un audit de conformité}
\begin{enumerate}
    \item Obtenir le registre des obligations, ou le construire avec la
          partie responsable.
    \item Communiquer les critères à la partie responsable avant le
          terrain.
    \item Fixer et écrire les seuils de signification, quantitatif et
          qualitatif.
    \item Coter chaque obligation sur chaque élément testé.
    \item Choisir la forme de la conclusion à partir de la lecture et des
          seuils.
\end{enumerate}
\end{boxmetier}

\begin{boxinfo}
\textbf{Ce que dit la recherche.} La cotation \NEK{} donne à la limitation
d'étendue une mesure~: le taux $U$~\cite{minka2026nek}. Elle permet de
distinguer, dans une conclusion avec réserve, ce qui relève d'un écart
établi (\NON{}) et ce qui relève d'une pièce manquante (\NEK{}), que les
rapports de conformité mêlent souvent sous le seul mot « réserve ».
\end{boxinfo}

\separateur

\begin{boxresume}
\textbf{Ce qu'il faut retenir de ce chapitre~:}
\begin{itemize}
    \item \textbf{L'audit de conformité} évalue le respect des critères
          applicables, dans une relation à trois parties, avec une assurance
          raisonnable ou limitée.
    \item \textbf{La méthode de l'ISSAI~4000}~: planification (critères,
          risques, seuils), exécution (preuves), rapport (conclusion,
          réponses, recommandations).
    \item \textbf{Le registre des obligations} écrit chaque obligation de
          façon atomique, avec son propriétaire, son contrôle et sa preuve.
    \item \textbf{La conclusion} est sans réserve, avec réserve,
          défavorable, ou impossible~; les seuils, quantitatif et
          qualitatif, la fixent.
    \item \textbf{\NEK{} mesure la limitation d'étendue}~; on ne projette
          pas un échantillon sans méthode.
    \item \textbf{Le droit des données personnelles} s'audite par thèmes, à
          partir du texte national vérifié.
\end{itemize}
\end{boxresume}

\begin{boxinfo}
\textbf{Points de contrôle de l'élève.} Je sais~: décrire la relation à
trois parties~; distinguer assurance raisonnable et limitée~; construire
une ligne du registre des obligations~; fixer un seuil qualitatif~;
choisir la forme d'une conclusion~; lister les thèmes d'un audit de
protection des données.
\end{boxinfo}

\bigskip

\begin{boxexo}[title={\ding{45}\quad Exercice~25.1. Raisonnable ou limitée~? \niveaubase}]
Pour chaque cas, choisissez le niveau d'assurance et rédigez la première
phrase de la conclusion~: (a) la Chambre des comptes teste 40~paiements et
rapproche toute la paie du fichier des départs~; (b) un consultant
interroge trois responsables et lit les procédures en deux jours.
\end{boxexo}

\begin{boxexo}[title={\ding{45}\quad Exercice~25.2. Le registre \niveaubase}]
Ajoutez au registre de l'ANPP deux obligations sur la conservation des
pièces de dépense et sur la sécurité du système de paiement, avec leur
contrôle et leur preuve.
\end{boxexo}

\begin{boxexo}[title={\ding{45}\quad Exercice~25.3. Le seuil \niveaumoyen}]
Le total des 40~paiements testés vaut 186~millions de FCFA. Calculez le
seuil quantitatif retenu par M.~NJOYA. Les 1~424~000~FCFA payés à des agents
partis le dépassent-ils~? Pourquoi la réserve est-elle quand même
nécessaire~?
\end{boxexo}

\begin{boxexo}[title={\ding{45}\quad Exercice~25.4. Recalculer la lecture \niveaumoyen}]
Recalculez $\mu$, $U$ et $H$ des paiements de l'ANPP (29/6/5). Les cinq
dossiers « aux archives » sont finalement produits~: trois sont conformes,
deux concernent encore des agents partis. Recalculez et dites si la forme
de la conclusion change.
\end{boxexo}

\begin{boxexo}[title={\ding{45}\quad Exercice~25.5. La forme de la conclusion \niveaumoyen}]
Choisissez la forme de la conclusion pour~: (a) 38/0/2~; (b) 20/18/2~; (c)
24/2/14~; (d) 36/1/3, où le \NON{} est un paiement à un bénéficiaire
fictif. Justifiez par le quadrant et les seuils.
\end{boxexo}

\begin{boxexo}[title={\ding{45}\quad Exercice~25.6. Projeter ou non \niveaumoyen}]
Le Parlement demande « combien l'ANPP a payé à tort en 2025 ». Que peut
répondre M.~NJOYA avec l'échantillon~? Quel test recommande-t-il pour
répondre précisément~?
\end{boxexo}

\begin{boxexo}[title={\ding{45}\quad Exercice~25.7. La fonction conformité \niveaubase}]
Distinguez, pour la BANQUE SANAGA, ce que fait le responsable de la
conformité et ce que fait l'audit interne à propos du règlement COBAC
R-2016/04.
\end{boxexo}

\begin{boxexo}[title={\ding{45}\quad Exercice~25.8. Le registre des traitements \niveaumoyen}]
Proposez les colonnes du registre des traitements de LiZbiZ et remplissez
une ligne pour le traitement « suivi des expéditions ».
\end{boxexo}

\begin{boxexo}[title={\ding{45}\quad Exercice~25.9. L'observation de l'ANPP \niveauexpert}]
Rédigez l'observation en sept éléments relative aux paiements versés à des
agents partis, avec sa gravité et ses recommandations, dont le test sur
toute la population.
\end{boxexo}

\begin{boxexo}[title={\ding{45}\quad Exercice~25.10. Un audit des données personnelles \niveauexpert}]
Préparez la lettre de mission d'un audit de conformité de LiZbiZ à la loi
n°~2024/017~: objectifs, critères (avec la mention de leur vérification
dans le texte officiel), niveau d'assurance, seuils, périmètre, et ce que
l'audit n'est pas.
\end{boxexo}

\bigskip

\begin{boxinfo}
\textbf{Pour aller plus loin.}
\begin{itemize}
    \item IDI, manuel d'audit de conformité~\cite{idi2021conformite}~: la
          méthode complète de l'ISSAI~4000.
    \item INTOSAI, GUID~5100~\cite{intosai2019guid5100}~: l'audit des SI
          dans le secteur public.
    \item ISO~37301:2021~\cite{iso37301_2021}~: le management de la
          conformité.
\end{itemize}
\end{boxinfo}

\begin{boxinfo}
\textbf{Transition.} La Partie~\ref{part:pratique} a audité chaque domaine
du SI. Le Chapitre~\ref{chap:cas-integrateurs} réunit les trois cas du
cours dans des missions complètes, de la lettre de mission au suivi.
\end{boxinfo}
